\cleardoublepage

\BookEntry[马特·里德利《基因组》]{马特·里德利《基因组》}

\BookCover{images/书单07基因组.jpg}{[英] 马特·里德利著| 李南哲译| 机械工业出版社2015年版}

\begin{BookSummaryBox}
\textbf{这本书全面介绍了人类基因组的发现对人类的重大意义}
\par
\end{BookSummaryBox}

\subsubsection*{这本书谈什么？}

人类基因组23对染色体的图谱绘制或许是21世纪迄今为止最重大的科学发现。它重新塑造了我们对疾病、寿命和自由意志的理解。
\par

《基因组》这本书全面解析了这一重大发现可能会有的深远影响。全书分23章，每一章会在23对染色体中的一对上挑选一个新近发现的基因，并讲述其故事。通过这种方式，作者介绍了人类这个物种从生命出现之初到未来生物学前沿的历史。书中探讨了由于基因组的图谱绘制而出现的科学、哲学以及伦理议题，这些将帮助我们每个人理解这一里程碑式的科学发现对人类意义。
\par

\subsubsection*{芒格为什么谈起这本书？}

在伯克希尔·哈撒韦2001财年的年会上，芒格说：
\par

\begin{quotation}
《基因组》分23章讲述了人类这个物种的历史。这绝对是一本令人惊叹的书，而且十分有趣。
\par
\end{quotation}

\subsubsection*{《基因组》作者马特·里德利如是说}

“人类基因组本身就是一部人类的自传，它从生命诞生之时起，便用‘基因语言’记录了人类和人类祖先所经历的世事更迭与变迁。”
\par

“读《基因组》，意味着我们能够更好地了解人类的起源、人类的进化、人类的本性，以及人类的思维，这将超过迄今为止科学带给我们的关于人类的所有知识。它将给人类学、心理学、医学、古生物学和几乎所有其他学科带来一场革命。”
\par

“基因组可以看作巨型的文档、巨大的图书、超长的菜谱，其所有内容都收录在比针尖还小的细胞的细胞核内。”
\par
